\section{Προσομοιωτής Αεροδυναμικής σε Python (AeroSandbox)}
\label{sec:impl_worker}

Ο worker προσομοίωσης υλοποιήθηκε σε Python με αξιοποίηση της βιβλιοθήκης AeroSandbox~\cite{sharpe2021aerosandbox}. Σε αντίθεση με το βαρύ και συχνά ασταθές σε σφάλματα πλεγματοποίησης OpenFOAM, το AeroSandbox προσφέρει ταχύτατη (\textasciitilde{}100\,ms) και ντετερμινιστική επίλυση 3D Vortex Lattice Method (VLM).

\subsection{Διατύπωση Μεθόδου Δικτύου Δινών (VLM)}
\label{subsec:impl_vlm}

Η μέθοδος VLM προσομοιώνει τη ροή γύρω από την πτέρυγα διακριτοποιώντας την επιφάνειά της σε $M \times N$ πλακίδια (panels). Σε κάθε πλακίδιο $(i,j)$ τοποθετείται ένα δεσμευμένο τμήμα δίνης (bound vortex segment) κατά το $1/4$ της χορδής. Η επαγόμενη ταχύτητα $\mathbf{u}_{ij}$ από κάθε τμήμα δίνης έντασης $\Gamma_{ij}$ υπολογίζεται μέσω του νόμου Biot-Savart:
\begin{equation}
\label{eq:biot_savart}
\mathbf{u} = \frac{\Gamma}{4\pi} \int_{\mathcal{C}} \frac{d\boldsymbol{\ell} \times \hat{\mathbf{r}}}{|\mathbf{r}|^2},
\end{equation}
όπου $\hat{\mathbf{r}}$ είναι το μοναδιαίο διάνυσμα από το στοιχείο $d\boldsymbol{\ell}$ προς το σημείο ελέγχου. Η συνθήκη μη-διαπερατότητας επιβάλλεται σε κάθε σημείο ελέγχου (collocation point) στο $3/4$ της χορδής, σχηματίζοντας γραμμικό σύστημα $\mathbf{A}\boldsymbol{\Gamma} = \mathbf{b}$ που επιλύεται με αποσύνθεση LU. Ο συντελεστής άνωσης πτέρυγας υπολογίζεται ως:
\begin{equation}
\label{eq:cl_wing}
C_L = \frac{2}{V_\infty S} \sum_{i,j} \Gamma_{ij} \, \Delta s_{ij},
\end{equation}
όπου $S$ η επιφάνεια αναφοράς πτέρυγας, $V_\infty$ η ταχύτητα ελεύθερης ροής και $\Delta s_{ij}$ το μήκος του $j$-οστού τμήματος δίνης στη σειρά $i$.

\subsection{Παραμετρικό Μοντέλο Αεροτομής NACA 4 Ψηφίων}
\label{subsec:impl_naca}

Η γεωμετρία της αεροτομής ορίζεται από την παραμετρική οικογένεια NACA 4 ψηφίων~\cite{abbott1959theory}. Η γραμμή καμπυλότητας (camber line) δίνεται από:
\begin{equation}
\label{eq:naca_camber}
y_c(x) =
\begin{cases}
\dfrac{m}{p^2}\bigl(2p x - x^2\bigr), & 0 \le x \le p, \\[6pt]
\dfrac{m}{(1-p)^2}\bigl(1 - 2p + 2p x - x^2\bigr), & p < x \le 1,
\end{cases}
\end{equation}
όπου $m$ είναι η μέγιστη καμπυλότητα (max camber) εκπεφρασμένη ως ποσοστό χορδής, $p$ η θέση μέγιστης καμπυλότητας ως δεκαδικό κλάσμα χορδής, και $x \in [0,1]$ η κανονικοποιημένη συντεταγμένη κατά μήκος χορδής. Η κατανομή πάχους προστίθεται συμμετρικά επί της γραμμής καμπυλότητας. Οι τρεις αεροδυναμικές παράμετροι ($m$, $p$, πάχος $t$) αποτελούν γονίδια του βελτιστοποιητή.

\subsection{Συνάρτηση Καταλληλότητας με Ποινές}
\label{subsec:impl_fitness}

Η συνάρτηση καταλληλότητας $F$ συνδυάζει τον αεροδυναμικό λόγο άνωσης-οπισθέλκουσας με ποινές παραβίασης περιορισμών:
\begin{equation}
\label{eq:fitness}
F = \frac{L}{D} - \lambda_1 \max\!\left(0,\; V_{\min} - V_{\text{fuse}}\right) - \lambda_2 \max\!\left(0,\; C_{m_\alpha}\right),
\end{equation}
όπου $L/D$ ο υπολογισμένος αεροδυναμικός βαθμός απόδοσης, $V_{\text{fuse}}$ ο διαθέσιμος εσωτερικός όγκος ατράκτου (σε mm$^3$), $V_{\min} = 20{.}000\,\text{mm}^3$ το ελάχιστο αποδεκτό όριο, και $C_{m_\alpha}$ η κλίση ροπής πρόνευσης (θετική τιμή υποδηλώνει αστάθεια. Οι συντελεστές ποινής ορίστηκαν ως $\lambda_1 = 10^{-3}$ ανά mm$^3$ παραβίασης και $\lambda_2 = 5$ ανά μονάδα $C_{m_\alpha}$, μετά από ρύθμιση με πειράματα ευαισθησίας.

\subsection{Χώρος Γονιδίων και Όρια Σχεδιασμού}
\label{subsec:impl_gene_bounds}

Το διάνυσμα σχεδιασμού αποτελείται από $d = 10$ γονίδια με φυσικά ερμηνεύσιμα όρια. Τα γονίδια διακρίνονται σε τρεις κατηγορίες: παράμετροι πτέρυγας (άνοιγμα $b \in [1.0, 4.0]$\,m, χορδή ρίζας $c_{\text{root}} \in [0.2, 0.8]$\,m, χορδή ακροπτερυγίου $c_{\text{tip}} \in [0.05, 0.4]$\,m, γωνία βέλους $\Lambda \in [0°, 45°]$, δίεδρο $\Gamma \in [-5°, 10°]$, συστροφή $\theta_{\text{twist}} \in [-5°, 5°]$), παράμετροι αεροτομής (καμπυλότητα $m \in [0, 0.09]$, θέση $p \in [0.1, 0.9]$, πάχος $t \in [0.06, 0.18]$) και κλίμακα ατράκτου (μήκος $L_{\text{fuse}} \in [0.5, 3.0]$\,m). Όλα τα γονίδια αναπαρίστανται ως συνεχείς πραγματικοί αριθμοί κανονικοποιημένοι σε $[0,1]$ κατά τη διάρκεια των τελεστών εξέλιξης.

\subsection{Παραγωγή Κλειδιού Hilbert για Αποθήκευση DHT}
\label{subsec:impl_hilbert_key}

Μετά την αξιολόγηση, κάθε άτομο αποθηκεύεται στο LEAD DHT με κλειδί Hilbert~\cite{skilling2004hilbert}. Το διάνυσμα γονιδίων $\mathbf{x} \in [0,1]^{10}$ κβαντίζεται σε $B = 16$-bit ακέραιες συντεταγμένες ανά διάσταση, και στη συνέχεια εφαρμόζεται ο αλγόριθμος κωδικοποίησης Hilbert για $d = 10$ διαστάσεις και $p = 16$ bits ανά διάσταση, παράγοντας ένα 64-bit κλειδί $h(\mathbf{x})$. Η ιδιότητα διατήρησης γειτνίασης (locality-preserving) της καμπύλης Hilbert εξασφαλίζει ότι γεωμετρικά παρόμοια σχέδια αντιστοιχούν σε κοντινά κλειδιά DHT, επιτρέποντας αποδοτικά ερωτήματα ομοιότητας.

\subsection{Συμβόλαιο gRPC: EvalRequest/EvalResponse}
\label{subsec:impl_grpc_worker}

Η επικοινωνία μεταξύ Orchestrator και εργατών γίνεται μέσω gRPC με μονοδιευθυντικές κλήσεις (unary RPC). Το μήνυμα \texttt{EvalRequest} περιέχει το 64-bit αναγνωριστικό ατόμου (\texttt{individual\_id}), το διάνυσμα γονιδίων (\texttt{genes: repeated double}), τον αριθμό γενεάς (\texttt{generation}) και το ανώτατο χρονικό όριο (\texttt{timeout\_ms}). Το μήνυμα \texttt{EvalResponse} επιστρέφει την τιμή καταλληλότητας (\texttt{fitness: double}), τον κωδικό κατάστασης (\texttt{status: EvalStatus}), τις επιμέρους αεροδυναμικές τιμές ($C_L$, $C_D$, $C_{m_\alpha}$, $V_{\text{fuse}}$) και το κλειδί Hilbert (\texttt{hilbert\_key: uint64}) έτοιμο για εισαγωγή στο DHT χωρίς επανυπολογισμό.

